
\documentclass{article}
\usepackage{colortbl}
\usepackage{makecell}
\usepackage{multirow}
\usepackage{supertabular}

\begin{document}

\newcounter{utterance}

\centering \large Interaction Transcript for game `dond', experiment `semi\_en', episode 4 with Aya{-}Expanse{-}32B{-}t0.0.
\vspace{24pt}

{ \footnotesize  \setcounter{utterance}{1}
\setlength{\tabcolsep}{0pt}
\begin{supertabular}{c@{$\;$}|p{.15\linewidth}@{}p{.15\linewidth}p{.15\linewidth}p{.15\linewidth}p{.15\linewidth}p{.15\linewidth}}
    \# & $\;$A & \multicolumn{4}{c}{Game Master} & $\;\:$B\\
    \hline

    \theutterance \stepcounter{utterance}  
    & & \multicolumn{4}{p{0.6\linewidth}}{
        \cellcolor[rgb]{0.9,0.9,0.9}{
            \makecell[{{p{\linewidth}}}]{
                \texttt{\tiny{[P1$\langle$GM]}}
                \texttt{You are playing a negotiation game in which you have to agree on how to divide a set of items among you and another player.} \\
\\ 
\texttt{Rules:} \\
\texttt{(a) You and the other player are given a set of items. Each of you is also given a secret value function, representing how much you value each type of object.} \\
\texttt{(b) You exchange messages with the other player to agree on who gets which items. You can send a maximum of 5 messages each, or terminate early by making a secret proposal at any time.} \\
\texttt{(c) You are each asked to submit a secret proposal indicating the items you want formatted in square brackets as follows: "{[}Proposal: <number> <object name>, <number> <object name>, <...>{]}"} \\
\texttt{(d) If your proposals are complementary, i.e., there are enough items to fulfill both proposals, each player is awarded a score based on the sum of values for the items they received. Otherwise, both of you get zero points.} \\
\texttt{(e) Your goal is to maximize the score you receive.} \\
\\ 
\texttt{Let us start.} \\
\\ 
\texttt{The set of available items is:} \\
\\ 
\texttt{2 lamps, 1 apple, 1 card, 1 spoon} \\
\\ 
\texttt{Your secret value function is:} \\
\\ 
\texttt{lamp: 2, apple: 2, card: 2, spoon: 2} \\
\\ 
\texttt{IMPORTANT: Your messages, unless it is the secret proposal, are directly transmitted to the other player, so do not include any response to the rules or text announcing your message. To make a secret proposal, use the indicated format. Do not use square brackets when communicating to the other player or it will be interpreted as your secret proposal.} \\
            }
        }
    }
    & \\ \\

    \theutterance \stepcounter{utterance}  
    & \multicolumn{4}{p{0.6\linewidth}}{
        \cellcolor[rgb]{0.9,0.9,0.9}{
            \makecell[{{p{\linewidth}}}]{
                \texttt{\tiny{[P1$\rangle$GM]}}
                \texttt{Understanding the rules and my valuation, I'll initiate the negotiation by messaging the other player:} \\
\\ 
\texttt{"Hello, I value all items equally. How about we split them evenly, assuming we can both agree on the values? I propose 1 lamp, 1 apple, and 1 card for me, and the remaining items for you."} \\
\\ 
\texttt{This message sets out a fair starting point given my valuation, and it leaves room for further negotiation based on the other player's response. If they agree to this split, we can both submit our secret proposals reflecting this agreement. If not, I will adjust my strategy based on their counteroffer or subsequent messages.} \\
\\ 
\texttt{Remember, my secret proposal (to be submitted later if necessary) would look like this:} \\
\texttt{"Proposal: 1 lamp, 1 apple, 1 card"} \\
            }
        }
    }
    & & \\ \\

    \theutterance \stepcounter{utterance}  
    & & & \multicolumn{2}{p{0.3\linewidth}}{
        \cellcolor[rgb]{0.9,0.9,0.9}{
            \makecell[{{p{\linewidth}}}]{
                \texttt{\tiny{[GM$|$GM]}}
                \texttt{continue} \\
            }
        }
    }
    & & \\ \\

    \theutterance \stepcounter{utterance}  
    & & & \multicolumn{2}{p{0.3\linewidth}}{
        \cellcolor[rgb]{0.9,0.9,0.9}{
            \makecell[{{p{\linewidth}}}]{
                \texttt{\tiny{[GM$|$GM]}}
                \texttt{Understanding the rules and my valuation, I'll initiate the negotiation by messaging the other player:} \\
\\ 
\texttt{"Hello, I value all items equally. How about we split them evenly, assuming we can both agree on the values? I propose 1 lamp, 1 apple, and 1 card for me, and the remaining items for you."} \\
\\ 
\texttt{This message sets out a fair starting point given my valuation, and it leaves room for further negotiation based on the other player's response. If they agree to this split, we can both submit our secret proposals reflecting this agreement. If not, I will adjust my strategy based on their counteroffer or subsequent messages.} \\
\\ 
\texttt{Remember, my secret proposal (to be submitted later if necessary) would look like this:} \\
\texttt{"Proposal: 1 lamp, 1 apple, 1 card"} \\
            }
        }
    }
    & & \\ \\

    \theutterance \stepcounter{utterance}  
    & & \multicolumn{4}{p{0.6\linewidth}}{
        \cellcolor[rgb]{0.9,0.9,0.9}{
            \makecell[{{p{\linewidth}}}]{
                \texttt{\tiny{[GM$\rangle$P2]}}
                \texttt{You are playing a negotiation game in which you have to agree on how to divide a set of items among you and another player.} \\
\\ 
\texttt{Rules:} \\
\texttt{(a) You and the other player are given a set of items. Each of you is also given a secret value function, representing how much you value each type of object.} \\
\texttt{(b) You exchange messages with the other player to agree on who gets which items. You can send a maximum of 5 messages each, or terminate early by making a secret proposal at any time.} \\
\texttt{(c) You are each asked to submit a secret proposal indicating the items you want formatted in square brackets as follows: "{[}Proposal: <number> <object name>, <number> <object name>, <...>{]}"} \\
\texttt{(d) If your proposals are complementary, i.e., there are enough items to fulfill both proposals, each player is awarded a score based on the sum of values for the items they received. Otherwise, both of you get zero points.} \\
\texttt{(e) Your goal is to maximize the score you receive.} \\
\\ 
\texttt{Let us start.} \\
\\ 
\texttt{The set of available items is:} \\
\\ 
\texttt{2 lamps, 1 apple, 1 card, 1 spoon} \\
\\ 
\texttt{Your secret value function is:} \\
\\ 
\texttt{lamp: 0, apple: 2, card: 3, spoon: 5} \\
\\ 
\texttt{IMPORTANT: Your messages, unless it is the secret proposal, are directly transmitted to the other player, so do not include any response to the rules or text announcing your message. To make a secret proposal, use the indicated format. Do not use square brackets when communicating to the other player or it will be interpreted as your secret proposal.} \\
\\ 
\\ 
\texttt{Understanding the rules and my valuation, I'll initiate the negotiation by messaging the other player:} \\
\\ 
\texttt{"Hello, I value all items equally. How about we split them evenly, assuming we can both agree on the values? I propose 1 lamp, 1 apple, and 1 card for me, and the remaining items for you."} \\
\\ 
\texttt{This message sets out a fair starting point given my valuation, and it leaves room for further negotiation based on the other player's response. If they agree to this split, we can both submit our secret proposals reflecting this agreement. If not, I will adjust my strategy based on their counteroffer or subsequent messages.} \\
\\ 
\texttt{Remember, my secret proposal (to be submitted later if necessary) would look like this:} \\
\texttt{"Proposal: 1 lamp, 1 apple, 1 card"} \\
            }
        }
    }
    & \\ \\

    \theutterance \stepcounter{utterance}  
    & & & \multicolumn{4}{p{0.6\linewidth}}{
        \cellcolor[rgb]{0.9,0.9,0.9}{
            \makecell[{{p{\linewidth}}}]{
                \texttt{\tiny{[GM$\langle$P2]}}
                \texttt{Understood. Here's my response to your initial proposal:} \\
\\ 
\texttt{"I appreciate your fair starting point. However, given my own valuation of the items, I must consider the spoon to be worth more than the combined value of the lamp, apple, and card. I'm willing to discuss a split that includes the spoon for me and the rest for you, with potential adjustments to ensure we both feel satisfied."} \\
\\ 
\texttt{This response indicates my interest in the spoon, while also leaving room for further negotiation. Based on your valuation, I expect you might also value the card highly, so we can explore options that benefit us both.} \\
\\ 
\texttt{Remember, I value the items as follows: lamp (0), apple (2), card (3), spoon (5). My secret proposal, if necessary, would reflect this valuation.} \\
            }
        }
    }
     \\ \\

    \theutterance \stepcounter{utterance}  
    & & & \multicolumn{2}{p{0.3\linewidth}}{
        \cellcolor[rgb]{0.9,0.9,0.9}{
            \makecell[{{p{\linewidth}}}]{
                \texttt{\tiny{[GM$|$GM]}}
                \texttt{continue} \\
            }
        }
    }
    & & \\ \\

    \theutterance \stepcounter{utterance}  
    & & & \multicolumn{2}{p{0.3\linewidth}}{
        \cellcolor[rgb]{0.9,0.9,0.9}{
            \makecell[{{p{\linewidth}}}]{
                \texttt{\tiny{[GM$|$GM]}}
                \texttt{Understood. Here's my response to your initial proposal:} \\
\\ 
\texttt{"I appreciate your fair starting point. However, given my own valuation of the items, I must consider the spoon to be worth more than the combined value of the lamp, apple, and card. I'm willing to discuss a split that includes the spoon for me and the rest for you, with potential adjustments to ensure we both feel satisfied."} \\
\\ 
\texttt{This response indicates my interest in the spoon, while also leaving room for further negotiation. Based on your valuation, I expect you might also value the card highly, so we can explore options that benefit us both.} \\
\\ 
\texttt{Remember, I value the items as follows: lamp (0), apple (2), card (3), spoon (5). My secret proposal, if necessary, would reflect this valuation.} \\
            }
        }
    }
    & & \\ \\

    \theutterance \stepcounter{utterance}  
    & & \multicolumn{4}{p{0.6\linewidth}}{
        \cellcolor[rgb]{0.9,0.9,0.9}{
            \makecell[{{p{\linewidth}}}]{
                \texttt{\tiny{[P1$\langle$GM]}}
                \texttt{Understood. Here's my response to your initial proposal:} \\
\\ 
\texttt{"I appreciate your fair starting point. However, given my own valuation of the items, I must consider the spoon to be worth more than the combined value of the lamp, apple, and card. I'm willing to discuss a split that includes the spoon for me and the rest for you, with potential adjustments to ensure we both feel satisfied."} \\
\\ 
\texttt{This response indicates my interest in the spoon, while also leaving room for further negotiation. Based on your valuation, I expect you might also value the card highly, so we can explore options that benefit us both.} \\
\\ 
\texttt{Remember, I value the items as follows: lamp (0), apple (2), card (3), spoon (5). My secret proposal, if necessary, would reflect this valuation.} \\
            }
        }
    }
    & \\ \\

    \theutterance \stepcounter{utterance}  
    & \multicolumn{4}{p{0.6\linewidth}}{
        \cellcolor[rgb]{0.9,0.9,0.9}{
            \makecell[{{p{\linewidth}}}]{
                \texttt{\tiny{[P1$\rangle$GM]}}
                \texttt{Thank you for sharing your valuation. It seems we have a clear understanding of each other's preferences. Given your interest in the spoon and my valuation, I'm willing to adjust my proposal to accommodate a fair split.} \\
\\ 
\texttt{Here's my counter{-}proposal:} \\
\\ 
\texttt{"I understand the spoon is valuable to you. How about this: I'll take 1 apple and 1 card, and you can have the spoon, 1 lamp, and any remaining item (if there is one). This way, we both get items we value highly."} \\
\\ 
\texttt{This proposal ensures that I receive items with a total value of 4 (apple + card) and you receive items worth 7 (spoon + lamp + potential second item), which seems fair given our valuations.} \\
\\ 
\texttt{If you agree, we can submit our secret proposals accordingly. If not, I'm open to further discussion to find a mutually beneficial solution.} \\
            }
        }
    }
    & & \\ \\

    \theutterance \stepcounter{utterance}  
    & & & \multicolumn{2}{p{0.3\linewidth}}{
        \cellcolor[rgb]{0.9,0.9,0.9}{
            \makecell[{{p{\linewidth}}}]{
                \texttt{\tiny{[GM$|$GM]}}
                \texttt{continue} \\
            }
        }
    }
    & & \\ \\

    \theutterance \stepcounter{utterance}  
    & & & \multicolumn{2}{p{0.3\linewidth}}{
        \cellcolor[rgb]{0.9,0.9,0.9}{
            \makecell[{{p{\linewidth}}}]{
                \texttt{\tiny{[GM$|$GM]}}
                \texttt{Thank you for sharing your valuation. It seems we have a clear understanding of each other's preferences. Given your interest in the spoon and my valuation, I'm willing to adjust my proposal to accommodate a fair split.} \\
\\ 
\texttt{Here's my counter{-}proposal:} \\
\\ 
\texttt{"I understand the spoon is valuable to you. How about this: I'll take 1 apple and 1 card, and you can have the spoon, 1 lamp, and any remaining item (if there is one). This way, we both get items we value highly."} \\
\\ 
\texttt{This proposal ensures that I receive items with a total value of 4 (apple + card) and you receive items worth 7 (spoon + lamp + potential second item), which seems fair given our valuations.} \\
\\ 
\texttt{If you agree, we can submit our secret proposals accordingly. If not, I'm open to further discussion to find a mutually beneficial solution.} \\
            }
        }
    }
    & & \\ \\

    \theutterance \stepcounter{utterance}  
    & & \multicolumn{4}{p{0.6\linewidth}}{
        \cellcolor[rgb]{0.9,0.9,0.9}{
            \makecell[{{p{\linewidth}}}]{
                \texttt{\tiny{[GM$\rangle$P2]}}
                \texttt{Thank you for sharing your valuation. It seems we have a clear understanding of each other's preferences. Given your interest in the spoon and my valuation, I'm willing to adjust my proposal to accommodate a fair split.} \\
\\ 
\texttt{Here's my counter{-}proposal:} \\
\\ 
\texttt{"I understand the spoon is valuable to you. How about this: I'll take 1 apple and 1 card, and you can have the spoon, 1 lamp, and any remaining item (if there is one). This way, we both get items we value highly."} \\
\\ 
\texttt{This proposal ensures that I receive items with a total value of 4 (apple + card) and you receive items worth 7 (spoon + lamp + potential second item), which seems fair given our valuations.} \\
\\ 
\texttt{If you agree, we can submit our secret proposals accordingly. If not, I'm open to further discussion to find a mutually beneficial solution.} \\
            }
        }
    }
    & \\ \\

    \theutterance \stepcounter{utterance}  
    & & & \multicolumn{4}{p{0.6\linewidth}}{
        \cellcolor[rgb]{0.9,0.9,0.9}{
            \makecell[{{p{\linewidth}}}]{
                \texttt{\tiny{[GM$\langle$P2]}}
                \texttt{Your counter{-}proposal is thoughtful and aligns better with our respective valuations. I agree to the split you've suggested:} \\
\\ 
\texttt{You: 1 apple, 1 card} \\
\texttt{Me: 1 spoon, 1 lamp, 1 remaining item (if applicable)} \\
\\ 
\texttt{This division ensures that we both receive items we value highly. Given your valuation, this translates to a score of 4 for you and 7 for me, which is a fair outcome.} \\
\\ 
\texttt{Let's submit our secret proposals to finalize this agreement:} \\
\\ 
\texttt{My secret proposal: "Proposal: 1 spoon, 1 lamp, any remaining item"} \\
\\ 
\texttt{I assume you will submit: "Proposal: 1 apple, 1 card"} \\
\\ 
\texttt{This concludes our negotiation successfully.} \\
            }
        }
    }
     \\ \\

    \theutterance \stepcounter{utterance}  
    & & & \multicolumn{2}{p{0.3\linewidth}}{
        \cellcolor[rgb]{0.9,0.9,0.9}{
            \makecell[{{p{\linewidth}}}]{
                \texttt{\tiny{[GM$|$GM]}}
                \texttt{continue} \\
            }
        }
    }
    & & \\ \\

    \theutterance \stepcounter{utterance}  
    & & & \multicolumn{2}{p{0.3\linewidth}}{
        \cellcolor[rgb]{0.9,0.9,0.9}{
            \makecell[{{p{\linewidth}}}]{
                \texttt{\tiny{[GM$|$GM]}}
                \texttt{Your counter{-}proposal is thoughtful and aligns better with our respective valuations. I agree to the split you've suggested:} \\
\\ 
\texttt{You: 1 apple, 1 card} \\
\texttt{Me: 1 spoon, 1 lamp, 1 remaining item (if applicable)} \\
\\ 
\texttt{This division ensures that we both receive items we value highly. Given your valuation, this translates to a score of 4 for you and 7 for me, which is a fair outcome.} \\
\\ 
\texttt{Let's submit our secret proposals to finalize this agreement:} \\
\\ 
\texttt{My secret proposal: "Proposal: 1 spoon, 1 lamp, any remaining item"} \\
\\ 
\texttt{I assume you will submit: "Proposal: 1 apple, 1 card"} \\
\\ 
\texttt{This concludes our negotiation successfully.} \\
            }
        }
    }
    & & \\ \\

    \theutterance \stepcounter{utterance}  
    & & \multicolumn{4}{p{0.6\linewidth}}{
        \cellcolor[rgb]{0.9,0.9,0.9}{
            \makecell[{{p{\linewidth}}}]{
                \texttt{\tiny{[P1$\langle$GM]}}
                \texttt{Your counter{-}proposal is thoughtful and aligns better with our respective valuations. I agree to the split you've suggested:} \\
\\ 
\texttt{You: 1 apple, 1 card} \\
\texttt{Me: 1 spoon, 1 lamp, 1 remaining item (if applicable)} \\
\\ 
\texttt{This division ensures that we both receive items we value highly. Given your valuation, this translates to a score of 4 for you and 7 for me, which is a fair outcome.} \\
\\ 
\texttt{Let's submit our secret proposals to finalize this agreement:} \\
\\ 
\texttt{My secret proposal: "Proposal: 1 spoon, 1 lamp, any remaining item"} \\
\\ 
\texttt{I assume you will submit: "Proposal: 1 apple, 1 card"} \\
\\ 
\texttt{This concludes our negotiation successfully.} \\
            }
        }
    }
    & \\ \\

    \theutterance \stepcounter{utterance}  
    & \multicolumn{4}{p{0.6\linewidth}}{
        \cellcolor[rgb]{0.9,0.9,0.9}{
            \makecell[{{p{\linewidth}}}]{
                \texttt{\tiny{[P1$\rangle$GM]}}
                \texttt{Excellent, we have reached a mutually agreeable solution. Our secret proposals, as discussed, will be:} \\
\\ 
\texttt{Your proposal: "{[}Proposal: 1 spoon, 1 lamp, any remaining item{]}"} \\
\\ 
\texttt{My proposal: "{[}Proposal: 1 apple, 1 card{]}"} \\
\\ 
\texttt{This negotiation has been a fair and productive exchange, demonstrating the power of clear communication and willingness to understand each other's preferences. Thank you for your cooperative approach.} \\
            }
        }
    }
    & & \\ \\

    \theutterance \stepcounter{utterance}  
    & & & \multicolumn{2}{p{0.3\linewidth}}{
        \cellcolor[rgb]{0.9,0.9,0.9}{
            \makecell[{{p{\linewidth}}}]{
                \texttt{\tiny{[GM$|$GM]}}
                \texttt{proposal does not include number/name pairs} \\
            }
        }
    }
    & & \\ \\

\end{supertabular}
}

\end{document}
